\chapter{The Dirac Gate as a First Variation}
\label{chap:dirac-gate-first-variation}

The Dirac equation has already appeared as the minimal first-order generator of a relativistic particle. It
returns now in a different role. It is no longer only the equation that carries a spinor forward. It is the
first-order face of a variational machine: the gate through which the universal tensor is tested. The gamma
matrices still square to the metric, the covariant derivative still carries the gauge connection, and the
curvature still appears as a commutator. But the question has changed. The apparatus now asks which state makes
the activity stationary when every admissible test variation is brought against it.

\section{The First-Order Gate}

Write the gauged Dirac operator as
\[
  \mathcal D_A \psi =
  i\gamma^\mu(\partial_\mu + iqA_\mu)\psi - m\psi .
\]
This is the first-order gate. The derivative is covariant because a spinor compared at nearby points must be
parallel-transported through the gauge frame before the comparison means anything. The gammas are Clifford
generators because the square of the first-order gate must return the metric face of the second-order wave
operator. The mass term closes the local particle law. The operator is therefore not an arbitrary differential
symbol. It is the smallest first-order object that can see spin, metric, and gauge at once.

For a variational computation, however, \(\mathcal D_A\psi=0\) is too narrow a place to stop. The real
experimental object is a residual. The state may meet a source, a boundary, a tensor background, or a finite
apparatus constraint. The useful form is
\[
  R(\psi) = \mathcal D_A\psi + \mathcal T\psi - f,
\]
where \(\mathcal T\) is the tensor term carried by the construction and \(f\) is the load supplied by the
experiment. A free Dirac equation is the special case \(\mathcal T=0\) and \(f=0\). The general case is the
interface Book II needs: a spinor tested against the tensor the construction implies.

\section{Why Fréchet}

A directional derivative along one coordinate is not enough. The apparatus varies a state as a whole: field
value, boundary trace, gauge response, spinor component, finite coefficient. The correct derivative is therefore
Fréchet. It asks for the linear map \(D\mathcal A(\psi)\) that best accounts for the change in an activity
\(\mathcal A\) under every small admissible perturbation \(\eta\):
\[
  \mathcal A(\psi+\eta)
  =
  \mathcal A(\psi)
  + D\mathcal A(\psi)[\eta]
  + o(\|\eta\|).
\]
The first variation is the bracketed linear functional. It is not merely a slope. It is the whole apparatus'
response to a test variation. When that response vanishes for all admissible \(\eta\), the state is stationary.

The activity used here is the squared Sobolev size of the residual:
\[
  \mathcal A(\psi) = \frac12\|R(\psi)\|_{H^s}^{2}.
\]
The Sobolev norm matters because the residual is not only a pointwise mismatch. It includes derivative content,
boundary regularity, and the finite smoothness floor at which the experiment is allowed to read. A state that is
small in raw value but violent in derivative is not small to the apparatus. The Sobolev norm is the ruler that
charges both.

\section{The Lean Spine}

The companion file is \(\texttt{device/Measurement/WeakDiracGalerkin.lean}\). It does not import a continuum PDE
library. It gives the exact finite spine:

\begin{quote}\small
\begin{verbatim}
abbrev Vec := List Int

structure UniversalTensor where
  diracRows : List Vec
  tensorRows : List Vec
  sobolevRows : List Vec
  load : Vec
deriving Repr
\end{verbatim}
\end{quote}

The rows are finite Galerkin rows. The Dirac rows represent the first-order gate. The tensor rows represent the
zeroth/universal tensor term. The Sobolev rows represent the norm in which the residual is tested. The load is
the experiment's right-hand side. Every entry is an integer in the companion file so that the example is exact.
Nothing is rounded. Nothing is hidden in a numerical tolerance.

The raw residual is computed by applying the Dirac rows and tensor rows to the state and subtracting the load:

\begin{quote}\small
\begin{verbatim}
def rawTensorResidual
    (U : UniversalTensor) (state : Vec) : Vec :=
  sub
    (add (matrixApply U.diracRows state)
      (matrixApply U.tensorRows state))
    U.load
\end{verbatim}
\end{quote}

This is the finite version of
\[
  R(\psi) = \mathcal D_A\psi + \mathcal T\psi - f.
\]
The chapter does not claim that a list of integers is a spinor field on a manifold. It claims something more
useful for an experiment: once the basis, boundary conditions, and coefficients have been chosen, the weak-form
residual becomes an exact finite object. That object can be checked.
